\section*{Capítulo \ref{ch:g11:igneous-rocks} --- As rochas ígneas}

\begin{solution}{exo:g11:igneous-rocks:1}
Uma \omterm{def:g11:igneous-rocks:volcanic-plutonic}{rocha vulcânica} se formou a partir de um magma que se solidificou na
superfície, como lava; uma \omterm{def:g11:igneous-rocks:volcanic-plutonic}{rocha plutônica} cristalizou em profundidade,
dentro de outras rochas. Pares: basalto e \omterm{def:g11:ocean-ridges:gabbro}{gabro} (\omterm{def:g11:igneous-rocks:felsic-mafic}{máficos}), \omterm{def:g11:igneous-rocks:rock-names}{andesito} e
\omterm{def:g11:igneous-rocks:rock-names}{diorito} (\omterm{def:g11:igneous-rocks:felsic-mafic}{intermediários}), \omterm{def:g11:igneous-rocks:rock-names}{riolito} e granito (\omterm{def:g11:igneous-rocks:felsic-mafic}{félsicos}).
\end{solution}

\begin{solution}{exo:g11:igneous-rocks:2}
(1) Vítrea. (2) De grão grosso. (3) \omterm{def:g11:igneous-rocks:porphyritic}{Porfirítica} (\omterm{def:g11:igneous-rocks:porphyritic}{fenocristais} de olivina
numa matriz fina). (4) De grão fino.
\end{solution}

\begin{solution}{exo:g11:igneous-rocks:3}
\qty{48}{\percent}: \omterm{def:g11:igneous-rocks:felsic-mafic}{máfico}, \omterm{def:g11:ocean-ridges:gabbro}{gabro}. \qty{60}{\percent}: \omterm{def:g11:igneous-rocks:felsic-mafic}{intermediário},
\omterm{def:g11:igneous-rocks:rock-names}{diorito}. \qty{72}{\percent}: \omterm{def:g11:igneous-rocks:felsic-mafic}{félsico}, granito.
\end{solution}

\begin{solution}{exo:g11:igneous-rocks:4}
Um \omterm{def:g11:igneous-rocks:intrusion}{dique} é uma lâmina de rocha ígnea que corta as camadas das rochas em que
penetra; uma \omterm{def:g11:igneous-rocks:intrusion}{soleira} é uma lâmina paralela às camadas, entre elas; um \omterm{def:g11:igneous-rocks:intrusion}{plúton}
é uma grande \omterm{def:g11:igneous-rocks:intrusion}{intrusão} de qualquer forma, formada em profundidade. O \omterm{def:g11:igneous-rocks:intrusion}{dique}
vertical que sobe do \omterm{def:g11:igneous-rocks:intrusion}{plúton} alimentou o vulcão.
\end{solution}

\begin{solution}{exo:g11:igneous-rocks:5}
A olivina se forma primeiro; o feldspato potássico, o quartzo e às vezes a
muscovita por último. O plagioclásio começa rico em cálcio e fica cada vez
mais rico em sódio à medida que o magma esfria.
\end{solution}

\begin{solution}{exo:g11:igneous-rocks:6}
(1) A: \qty{74}{\percent} de \ce{SiO2}, \omterm{def:g11:igneous-rocks:felsic-mafic}{félsica}; B: \qty{47}{\percent},
\omterm{def:g11:igneous-rocks:felsic-mafic}{máfica}. (2) A é pobre em \ce{FeO}, \ce{MgO} e \ce{CaO} e rica em \ce{Na2O} e
\ce{K2O}, como um magma \omterm{def:g11:igneous-rocks:felsic-mafic}{félsico}; B é rica em \ce{FeO}, \ce{MgO} e \ce{CaO},
pobre em \ce{Na2O}, como um magma \omterm{def:g11:igneous-rocks:felsic-mafic}{máfico}. (3) A é um \omterm{def:g11:igneous-rocks:rock-names}{riolito}; B, um \omterm{def:g11:ocean-ridges:gabbro}{gabro}.
\end{solution}

\begin{solution}{exo:g11:igneous-rocks:7}
O tempo cresce como o quadrado da espessura. \qty{1}{m}: $10/4 = 2.5$
dias. \qty{20}{m}: $10 \times 10^2 = 1000$ dias, cerca de 2.7 anos.
\qty{200}{m}: $10 \times 100^2 = \num{100000}$ dias, cerca de 270 anos.
\end{solution}

\begin{solution}{exo:g11:igneous-rocks:8}
(1) \unit{m^2} dividido por \unit{m^2/s} dá segundos. (2) $400/10^{-6} =
\qty{4e8}{s}$, isto é, $4 \times 10^8 / 3.16 \times 10^7 \approx 12.7$
anos. (3) $0.21 \times 12.7 \approx 2.7$ anos.
\end{solution}

\begin{solution}{exo:g11:igneous-rocks:9}
Ela é \omterm{def:g11:igneous-rocks:porphyritic}{porfirítica}: \omterm{def:g11:igneous-rocks:porphyritic}{fenocristais} de plagioclásio numa matriz de grão fino.
Primeiro o magma esfriou devagar em profundidade, e os cristais de
plagioclásio cresceram muito enquanto a maior parte do magma continuava
líquida; depois ele subiu e esfriou depressa, e o líquido restante
cristalizou formando a matriz fina.
\end{solution}

\begin{solution}{exo:g11:igneous-rocks:10}
(1) À esquerda: 2 Mg, 2 Si, 6 O; à direita: 2 Mg, 2 Si, 6 O. (2) \qty{1}{kg}
de olivina são $1000/140.7 \approx 7.11$~mol, que tiram o mesmo número de
mols de sílica: $7.11 \times 60.1 \approx \qty{427}{g}$. (3) Sem sílica
restante no líquido, a reação não pode continuar: a olivina fica.
\end{solution}

\begin{solution}{exo:g11:igneous-rocks:11}
(1) $2.5 \times 20 \times 1100 = \qty{55000}{MJ}$, \qty{5.5e10}{J}.
(2) $900 \times 20 = \qty{18000}{MJ}$, \qty{1.8e10}{J}. (3) O calor de
cristalização é cerca de um terço do calor perdido pelo resfriamento, um
quarto do total: o bastante para retardar muito o congelamento, porque ele é
todo liberado acima do solidus.
\end{solution}

\begin{solution}{exo:g11:igneous-rocks:12}
(1) Os cristais menores no contato, que congelou em menos de um dia; os
maiores no centro, que continuou líquido por cerca de 950 dias. (2) Não: os
primeiros \qty{5}{m} levaram cerca de 545 dias, os últimos \qty{5}{m} cerca
de 400 dias; a frente desacelera no início, e depois acelera perto do
centro, onde o calor que sai pelos dois lados se encontra. (3) Mesmo no
centro, a \omterm{def:g11:igneous-rocks:intrusion}{soleira} congelou em menos de três anos; um granito de cristais
grandes precisa de um tempo muito mais longo.
\end{solution}

\begin{solution}{exo:g11:igneous-rocks:13}
(1) O \omterm{def:g11:igneous-rocks:intrusion}{plúton} é 100 vezes mais espesso: $2.6 \times 100^2 = \num{26000}$
anos. (2) Por exemplo: um \omterm{def:g11:igneous-rocks:intrusion}{plúton} profundo fica numa rocha que já está
quente, e por isso perde calor mais devagar do que para uma rocha a
\qty{50}{\degreeCelsius}; ele pode ser realimentado por magma novo várias
vezes; um \omterm{def:g11:igneous-rocks:intrusion}{plúton} não é uma lâmina infinita; as propriedades da rocha mudam
com a temperatura.
\end{solution}

\begin{solution}{exo:g11:igneous-rocks:14}
Um magma \omterm{def:g11:igneous-rocks:felsic-mafic}{máfico} é quente e pobre em sílica: ele cristaliza os primeiros
minerais da série --- olivina, piroxênio, plagioclásio rico em cálcio --- e
esgota a sua sílica antes que o quartzo possa se formar. Um magma \omterm{def:g11:igneous-rocks:felsic-mafic}{félsico} é
mais frio e rico em sílica: a olivina não é estável nele, e ele chega ao fim
da série, formando plagioclásio rico em sódio, feldspato potássico e
quartzo.
\end{solution}

\begin{solution}{exo:g11:igneous-rocks:15}
(1) $b = \sqrt{t\kappa/C} = \sqrt{86400 \times 10^{-6}/0.21} \approx
\qty{0.64}{m}$. (2) A lei dá $0.21 \times 25/10^{-6} \approx
\qty{5.3e6}{s}$, cerca de dois meses: a mesma ordem que um mês (um derrame
de lava também perde calor para o ar pelo seu topo). (3) Vítrea para o
material lançado no ar ou na água; de grão fino para os derrames de lava e
as lâminas finas; mais grossa, mas ainda pequena, para as \omterm{def:g11:igneous-rocks:intrusion}{intrusões} rasas e
espessas.
\end{solution}

\begin{solution}{pb:g11:igneous-rocks:1}
\textbf{1.} Uma \omterm{def:g11:igneous-rocks:intrusion}{soleira}: ela é paralela às camadas, entre elas.

\textbf{2.} No contato, de grão muito fino, quase vítrea; no centro, de grão
fino a médio, com os minerais mal visíveis com uma lupa.

\textbf{3.} O contato perdeu o seu calor mais depressa, para o arenito
frio: os cristais não tiveram tempo de crescer.

\textbf{4.} Uma margem resfriada. Ela prova que o magma estava quente quando
se instalou, e o arenito frio.

\textbf{5.} Um \omterm{def:g11:igneous-rocks:intrusion}{dique}, o alimentador da \omterm{def:g11:igneous-rocks:intrusion}{soleira}.

\textbf{6.} \omterm{def:g11:igneous-rocks:felsic-mafic}{Máfico}: \qty{49}{\percent} de \ce{SiO2}, rico em ferro, magnésio
e cálcio.

\textbf{7.} Olivina, piroxênio e plagioclásio rico em cálcio.

\textbf{8.} Um basalto; a \omterm{def:g11:igneous-rocks:volcanic-plutonic}{rocha plutônica} de mesma composição é o \omterm{def:g11:ocean-ridges:gabbro}{gabro}.

\textbf{9.} A \omterm{def:g11:igneous-rocks:texture}{textura} depende da velocidade de resfriamento, que difere do
contato ao centro, e não da composição.

\textbf{10.} A olivina primeiro; o plagioclásio começando mais ou menos no
mesmo ponto que o piroxênio.

\textbf{11.} $b^2/\kappa = 400/10^{-6} = \qty{4e8}{s}$.

\textbf{12.} $0.206 \times 4 \times 10^8 \approx \qty{8.2e7}{s}$, isto é,
$8.2 \times 10^7 / 3.16 \times 10^7 \approx 2.6$ anos.

\textbf{13.} $0.206 \times 4/10^{-6} \approx \qty{8.2e5}{s}$, cerca de 9.5
dias.

\textbf{14.} $0.07 \times 4 \times 10^8 = \qty{2.8e7}{s}$, cerca de 0.9 ano.
O magma que cristaliza libera calor, que também precisa ser levado embora
antes que o magma se torne sólido.

\textbf{15.} Cerca de 950 dias contra 0.6 dia: cerca de 1600 vezes mais depressa.

\textbf{16.} Cerca de 950 dias, uns 2.6 anos: muito mais que um derrame de
lava, dentro dos meses a anos das \omterm{def:g11:igneous-rocks:intrusion}{intrusões} rasas.

\textbf{17.} Mil vezes mais tempo exige uma espessura $\sqrt{1000} \approx
32$ vezes maior: cerca de \qty{630}{m}.

\textbf{18.} Perto da superfície, um corpo de magma irromperia, ou esfriaria
depressa demais; um corpo grande que esfria durante muito tempo precisa
estar no fundo, numa rocha que já está quente.

\textbf{19.} Elas foram desgastadas pela erosão.

\textbf{20.} A \omterm{def:g11:igneous-rocks:intrusion}{soleira} levou cerca de \textbf{2.6 anos} para se solidificar.
\end{solution}
